\documentclass[11pt]{article}
\usepackage[a4paper,margin=1in]{geometry}
\usepackage{enumitem}
\usepackage{hyperref}
\usepackage{amsmath}

\title{Mental Math Trainer --- Architecture}
\date{1 March 2026}

\begin{document}

\maketitle

\section{Purpose}

This document defines the architecture for a mental math trainer built with:

\begin{itemize}[leftmargin=*]
	\item Next.js App Router + TypeScript
	\item PostgreSQL + Prisma
	\item better-auth
	\item Tailwind CSS
\end{itemize}

Primary objective: support a domain-based difficulty system (ADD, SUB, MUL, DIV) with deterministic question generation and user progression.

\section{Current State (as implemented)}

\subsection*{App Layers}

\begin{itemize}[leftmargin=*]
	\item \textbf{UI layer} (App Router pages):
	\begin{itemize}
		\item Dashboard setup: choose operations and question count.
		\item Training screen: render questions, timer, answer input, skip flow.
		\item Overview page exists but is currently empty.
	\end{itemize}
	\item \textbf{API layer}:
	\begin{itemize}
		\item \texttt{GET /api/training/calculations} returns generated questions.
		\item Endpoint requires an authenticated session.
	\end{itemize}
	\item \textbf{Domain layer}:
	\begin{itemize}
		\item \texttt{question-generator.ts} parses pattern params and generates deterministic questions from \texttt{(domain, params, seed)}.
	\end{itemize}
	\item \textbf{Persistence layer} (Prisma):
	\begin{itemize}
		\item Pattern catalog exists.
		\item Attempt and user progress tables exist.
		\item Training flow currently does \emph{not} persist attempts yet.
	\end{itemize}
\end{itemize}

\subsection*{Important Current Limitation}

The question API currently fetches only \texttt{level = 1} patterns. That means difficulty progression data model exists, but progression logic is not yet integrated in the serving path.

\section{Core Domain Model}

\subsection*{Entities}

\begin{itemize}[leftmargin=*]
	\item \textbf{User}: authentication identity.
	\item \textbf{Pattern}: difficulty unit per domain and level; contains generation constraints in \texttt{params} JSON and a \texttt{cutoffTimeMs} target.
	\item \textbf{Attempt}: one user interaction with one generated question instance.
	\item \textbf{UserDomainProgress}: current and highest unlocked level per user and domain.
\end{itemize}

\subsection*{Question Identity}

A generated question is represented by:

\begin{center}
\texttt{(patternId, seed)}
\end{center}

With deterministic RNG, this pair is enough to reconstruct the prompt and expected answer.

\section{Request and Data Flow}

\subsection*{Read Flow (today)}

\begin{enumerate}[leftmargin=*]
	\item Client reads selected operations and count from URL query.
	\item Client calls \texttt{/api/training/calculations}.
	\item API verifies session via better-auth.
	\item API loads active patterns (currently fixed to level 1).
	\item API generates deterministic questions from seeds.
	\item API returns list to client for local session execution.
\end{enumerate}

\subsection*{Write Flow (intended architecture)}

\begin{enumerate}[leftmargin=*]
	\item Client submits result for each question: \texttt{patternId}, \texttt{seed}, answer metadata, timings.
	\item API validates payload and recomputes expected answer server-side.
	\item API writes \texttt{Attempt} row.
	\item API updates \texttt{UserDomainProgress} according to progression algorithm.
\end{enumerate}

\section{Difficulty System Architecture}

\subsection*{Design Principles}

\begin{itemize}[leftmargin=*]
	\item Progression is \textbf{per domain}, not global.
	\item A level maps to one or more \textbf{patterns}.
	\item Serving and progression decisions are made from persisted attempts.
	\item Question generation remains stateless and deterministic.
\end{itemize}

\subsection*{Serving Strategy}

For each selected domain, resolve the user's current level from \texttt{UserDomainProgress}, then sample patterns with:

\begin{itemize}[leftmargin=*]
	\item majority from current level $L$
	\item small slice from $L-1$ (stability)
	\item small slice from $L+1$ (challenge)
\end{itemize}

Exact percentages are a policy parameter and can be tuned independently of schema.

\subsection*{Promotion/Demotion Ownership}

Promotion/demotion logic must run server-side during attempt ingestion. The client should never decide level changes.

\section{Module Boundaries}

\begin{itemize}[leftmargin=*]
	\item \textbf{Auth module}: \texttt{lib/auth.ts}, \texttt{lib/session.ts}, auth API route.
	\item \textbf{Training generation module}: \texttt{features/training/question-generator.ts}.
	\item \textbf{Training fetch client module}: \texttt{features/training/use-training-questions.ts}.
	\item \textbf{Training runtime UI}: \texttt{app/training/page.tsx}.
	\item \textbf{Persistence module}: Prisma models and queries.
\end{itemize}

Rule: generation logic is pure; persistence and progression stay in API/server modules.

\section{Security and Integrity}

\begin{itemize}[leftmargin=*]
	\item All training APIs require session.
	\item Attempt correctness is computed server-side from \texttt{patternId + seed}; do not trust client-side \texttt{solved} flags.
	\item Keep pattern constraints authoritative in DB; client is presentation only.
\end{itemize}

\section{Incremental Delivery Plan}

\textbf{Step 1 (this doc):} architecture baseline and module boundaries.

\textbf{Step 2:} add attempt ingestion endpoint (POST), with server-side answer verification.

\textbf{Step 3:} wire training UI submit/skip actions to persist attempts.

\textbf{Step 4:} implement level resolver (read \texttt{UserDomainProgress}) and stop serving only level 1.

\textbf{Step 5:} implement progression updater (promotion/demotion policy).

\textbf{Step 6:} implement overview metrics from real attempts (replace mocks).


\end{document}
